% =====================================================================
\chapter{Problemas de optimización y convexidad}
\label{tema:1}
% =====================================================================

Este tema plantea el entrenamiento de un modelo como un problema de minimización sin restricciones y reúne las herramientas de análisis que se usan en el resto del módulo: convexidad, convexidad fuerte, suavidad de Lipschitz y condiciones de optimalidad. Las definiciones y los resultados siguen los capítulos~2, 3 y~9 de~\cite{bv}; la discusión sobre redes profundas sigue~\cite[§8.1 y §8.2]{gbc} y~\cite[§12.1 y §12.2]{d2l}. Parte de la materia de convexidad se introdujo en el módulo~4 (tema~5); aquí se completa con las constantes $\mu$ y $L$, que determinan la velocidad de los algoritmos.

\section{Minimización del riesgo empírico}

\begin{definicion}[Problema de optimización sin restricciones]
Sea $f\colon\R^{n}\to\R$. El problema \emph{minimizar $f(\theta)$, $\theta\in\R^{n}$}, consiste en hallar el \emph{valor óptimo} $f^{\ast}=\inf_{\theta}f(\theta)$ y, si existe, un \emph{minimizador global} $\theta^{\ast}$ con $f(\theta^{\ast})=f^{\ast}$. Un punto $\bar\theta$ es un \emph{minimizador local} si existe $r>0$ tal que $f(\bar\theta)\le f(\theta)$ para todo $\theta$ con $\norm{\theta-\bar\theta}<r$.
\end{definicion}

\begin{definicion}[Riesgo empírico]\label{def:riesgo}
Sean $N$ ejemplos de entrenamiento y $\ell_i(\theta)$ la pérdida del modelo de parámetros $\theta$ sobre el ejemplo $i$. El \emph{riesgo empírico} es
\[
J(\theta)=\frac1N\sum_{i=1}^{N}\ell_i(\theta).
\]
\end{definicion}

En un modelo de lenguaje, cada ejemplo es una posición de un texto y $\ell_i(\theta)=-\ln p_\theta(x_t\mid x_{1:t-1})$ es la log-verosimilitud negativa del token observado; $J$ es la pérdida $\mathcal{J}$ del módulo~1 (tema~5). El objetivo real no es minimizar $J$, sino el riesgo esperado sobre la distribución de los datos, al que $J$ solo aproxima; por eso el aprendizaje se diferencia de la optimización pura~\cite[§8.1]{gbc}. Este módulo trata, sin embargo, del problema de optimización.

\section{Convexidad, convexidad fuerte y suavidad}

\begin{definicion}[Conjunto y función convexos]
Un conjunto $C\subseteq\R^{n}$ es \emph{convexo} si $t x+(1-t)y\in C$ para todo $x,y\in C$ y $t\in[0,1]$. Una función $f\colon\R^{n}\to\R$ es \emph{convexa} si
\[
f\bigl(tx+(1-t)y\bigr)\le t f(x)+(1-t)f(y)\qquad\text{para todo } x,y\in\R^{n},\ t\in[0,1],
\]
y \emph{estrictamente convexa} si la desigualdad es estricta cuando $x\ne y$ y $t\in(0,1)$~\cite[cap.~2 y §3.1]{bv}.
\end{definicion}

\begin{teorema}[Caracterización de primer orden]\label{teo:primer-orden}
Sea $f$ diferenciable. Entonces $f$ es convexa si y solo si
\[
f(y)\ge f(x)+\nabla f(x)\T(y-x)\qquad\text{para todo } x,y\in\R^{n}.
\]
\end{teorema}

\begin{proof}
($\Rightarrow$) Sean $x,y$ y $t\in(0,1]$. Por convexidad, $f(x+t(y-x))\le f(x)+t\bigl(f(y)-f(x)\bigr)$, es decir,
\[
\frac{f(x+t(y-x))-f(x)}{t}\le f(y)-f(x).
\]
Cuando $t\to 0^{+}$, el primer miembro tiende a la derivada direccional $\nabla f(x)\T(y-x)$, lo que da la desigualdad.

($\Leftarrow$) Sean $x,y$, $t\in[0,1]$ y $z=tx+(1-t)y$. Aplicando la hipótesis en el punto $z$ a $x$ y a $y$,
$f(x)\ge f(z)+\nabla f(z)\T(x-z)$ y $f(y)\ge f(z)+\nabla f(z)\T(y-z)$. Multiplicando por $t$ y por $1-t$ y sumando, el término lineal es $\nabla f(z)\T\bigl(tx+(1-t)y-z\bigr)=0$, y queda $tf(x)+(1-t)f(y)\ge f(z)$.
\end{proof}

Geométricamente, una función convexa queda por encima de todas sus rectas (o hiperplanos) tangentes. Para funciones dos veces diferenciables, la convexidad equivale a $\nabla^{2}f(x)\succeq 0$ para todo $x$ (módulo~4, tema~5; \cite[§3.1]{bv}). Las dos definiciones siguientes refuerzan esta condición por abajo y la limitan por arriba.

\begin{definicion}[Convexidad fuerte]\label{def:fuerte}
Sea $\mu>0$. Una función diferenciable $f$ es \emph{$\mu$-fuertemente convexa} si
\[
f(y)\ge f(x)+\nabla f(x)\T(y-x)+\frac{\mu}{2}\norm{y-x}^{2}\qquad\text{para todo } x,y.
\]
Si $f$ es dos veces diferenciable, esto equivale a $\nabla^{2}f(x)\succeq\mu I$ para todo $x$~\cite[§9.1]{bv}.
\end{definicion}

\begin{definicion}[Suavidad de Lipschitz]\label{def:suave}
Sea $L>0$. Una función diferenciable $f$ es \emph{$L$-suave} si su gradiente es lipschitziano de constante $L$:
\[
\norm{\nabla f(x)-\nabla f(y)}\le L\norm{x-y}\qquad\text{para todo } x,y.
\]
Si $f$ es dos veces diferenciable, esto equivale a que todos los autovalores de $\nabla^{2}f(x)$ estén en $[-L,L]$ para todo $x$; para $f$ convexa, a $\nabla^{2}f(x)\preceq LI$.
\end{definicion}

Si $f$ es $\mu$-fuertemente convexa y $L$-suave, necesariamente $\mu\le L$, y el cociente
\[
\kappa=\frac{L}{\mu}\ge 1
\]
se llama \emph{número de condición} de $f$. En el \cref{tema:2} se demuestra que $\kappa$ controla el número de iteraciones del descenso de gradiente.

\begin{ejemplo}[Funciones cuadráticas]\label{ej:cuadratica}
Sea $f(\theta)=\tfrac12\theta\T Q\theta+c\T\theta+b$, con $Q\in\R^{n\times n}$ simétrica, $c\in\R^{n}$ y $b\in\R$. Entonces $\nabla f(\theta)=Q\theta+c$ y $\nabla^{2}f=Q$ es constante. Si $Q\succ 0$, la función es $\mu$-fuertemente convexa y $L$-suave con
\[
\mu=\lambda_{\min}(Q),\qquad L=\lambda_{\max}(Q),\qquad \kappa=\frac{\lambda_{\max}(Q)}{\lambda_{\min}(Q)},
\]
que es el número de condición espectral de $Q$ del módulo~3. Para $f(x,y)=x^{2}+5y^{2}$ se tiene $Q=\diag(2,10)$, $\mu=2$, $L=10$ y $\kappa=5$: las curvas de nivel son elipses cuyos semiejes están en razón $\sqrt\kappa=\sqrt5\approx 2{,}236$.
\end{ejemplo}

\section{Condiciones de optimalidad}

\begin{teorema}[Condición necesaria de primer orden]\label{teo:fermat}
Si $f$ es diferenciable y $\bar\theta$ es un minimizador local, entonces $\nabla f(\bar\theta)=0$.
\end{teorema}

\begin{proof}
Sea $d=-\nabla f(\bar\theta)$ y supóngase $d\ne 0$. La función $\varphi(t)=f(\bar\theta+td)$ cumple $\varphi'(0)=\nabla f(\bar\theta)\T d=-\norm{\nabla f(\bar\theta)}^{2}<0$, luego $\varphi(t)<\varphi(0)$ para todo $t>0$ suficientemente pequeño. Esto contradice que $\bar\theta$ sea minimizador local.
\end{proof}

Un punto con $\nabla f(\bar\theta)=0$ se llama \emph{punto crítico} o \emph{estacionario}. En general, la condición no es suficiente; para funciones convexas, sí lo es.

\begin{teorema}[Optimalidad en el caso convexo]\label{teo:opt-convexa}
Sea $f$ convexa y diferenciable. Un punto $\theta^{\ast}$ es minimizador global si y solo si $\nabla f(\theta^{\ast})=0$. En particular, todo minimizador local de una función convexa es global.
\end{teorema}

\begin{proof}
Si $\theta^{\ast}$ es minimizador global, también es local y el \cref{teo:fermat} da $\nabla f(\theta^{\ast})=0$. Recíprocamente, si $\nabla f(\theta^{\ast})=0$, el \cref{teo:primer-orden} da $f(y)\ge f(\theta^{\ast})+0$ para todo $y$. La última afirmación se sigue de combinar ambas implicaciones con el \cref{teo:fermat}.
\end{proof}

\begin{proposicion}[Minimizador de una cuadrática]\label{prop:min-cuadratica}
Sea $f(\theta)=\tfrac12\theta\T Q\theta+c\T\theta+b$ con $Q\succ 0$. Su único minimizador y su valor mínimo son
\[
\theta^{\ast}=-Q^{-1}c,\qquad f^{\ast}=b-\tfrac12\,c\T Q^{-1}c,
\]
y para todo $\theta$ se cumple $f(\theta)-f^{\ast}=\tfrac12(\theta-\theta^{\ast})\T Q(\theta-\theta^{\ast})$.
\end{proposicion}

\begin{proof}
$f$ es convexa porque $\nabla^{2}f=Q\succ 0$. Por el \cref{teo:opt-convexa}, los minimizadores son las soluciones de $Q\theta+c=0$; como $Q$ es invertible (sus autovalores son positivos), la única es $\theta^{\ast}=-Q^{-1}c$. Desarrollando $\tfrac12(\theta-\theta^{\ast})\T Q(\theta-\theta^{\ast})$ con $Q\theta^{\ast}=-c$ se obtiene $\tfrac12\theta\T Q\theta+c\T\theta+\tfrac12\theta^{\ast\top}Q\theta^{\ast}$, y $\tfrac12\theta^{\ast\top}Q\theta^{\ast}=\tfrac12c\T Q^{-1}c$. Por tanto $f(\theta)=\tfrac12(\theta-\theta^{\ast})\T Q(\theta-\theta^{\ast})+b-\tfrac12c\T Q^{-1}c$, que es mínimo, con valor $b-\tfrac12c\T Q^{-1}c$, cuando $\theta=\theta^{\ast}$.
\end{proof}

La identidad de la \cref{prop:min-cuadratica} permite trasladar el minimizador al origen: en lo que sigue, al estudiar cuadráticas se supondrá sin pérdida de generalidad $f(\theta)=\tfrac12\theta\T Q\theta$, con $\theta^{\ast}=0$ y $f^{\ast}=0$.

\begin{proposicion}[Desigualdad de Polyak--Łojasiewicz]\label{prop:pl}
Si $f$ es $\mu$-fuertemente convexa y diferenciable, con valor óptimo $f^{\ast}$, entonces
\[
\norm{\nabla f(x)}^{2}\ge 2\mu\bigl(f(x)-f^{\ast}\bigr)\qquad\text{para todo } x.
\]
\end{proposicion}

\begin{proof}
Fijado $x$, el segundo miembro de la \cref{def:fuerte} es una cuadrática en $y$ que se minimiza en $y=x-\nabla f(x)/\mu$, con valor $f(x)-\norm{\nabla f(x)}^{2}/(2\mu)$. Por tanto, $f(y)\ge f(x)-\norm{\nabla f(x)}^{2}/(2\mu)$ para todo $y$. Tomando el ínfimo en $y$, $f^{\ast}\ge f(x)-\norm{\nabla f(x)}^{2}/(2\mu)$, que es la desigualdad buscada.
\end{proof}

La \cref{prop:pl} dice que, lejos del óptimo en valor, el gradiente no puede ser pequeño. Es la propiedad que hace converger linealmente al descenso de gradiente (\cref{teo:gd-fuerte}).

\section{Paisajes de pérdida no convexos}

La pérdida de una red neuronal profunda no es convexa en sus parámetros (módulo~4, tema~5). La información local en un punto crítico la proporciona el hessiano.

\begin{proposicion}[Clasificación de puntos críticos]\label{prop:clasificacion}
Sea $f$ de clase $C^{2}$ y $\bar\theta$ un punto crítico.
\begin{enumerate}
  \item Si $\nabla^{2}f(\bar\theta)\succ 0$, $\bar\theta$ es un minimizador local estricto.
  \item Si $\nabla^{2}f(\bar\theta)$ tiene un autovalor negativo, $\bar\theta$ no es minimizador local; si tiene autovalores de ambos signos, es un \emph{punto de silla}.
  \item Si $\nabla^{2}f(\bar\theta)\succeq 0$ con algún autovalor nulo, el criterio no decide.
\end{enumerate}
\end{proposicion}

La demostración, basada en el desarrollo de Taylor de segundo orden, se dio en el módulo~4 (tema~4). El \cref{tab:criticos} recoge los casos con ejemplos en $\R^{2}$.

\begin{table}[h]
\centering\small
\caption{Puntos críticos en el origen y autovalores del hessiano.}\label{tab:criticos}
\begin{tabular}{@{}llll@{}}
\toprule
Tipo & Autovalores de $\nabla^{2}f$ & Ejemplo & Comportamiento del descenso \\
\midrule
Mínimo local & todos $>0$ & $x^{2}+y^{2}$ & converge si el paso es pequeño \\
Máximo local & todos $<0$ & $-x^{2}-y^{2}$ & se aleja \\
Punto de silla & de ambos signos & $x^{2}-y^{2}$ & se aleja salvo en un conjunto de medida nula \\
Valle plano & alguno nulo & $x^{2}$ & no hay fuerza en la dirección plana \\
\bottomrule
\end{tabular}
\end{table}

\begin{ejemplo}[Silla]\label{ej:silla}
Para $f(x,y)=x^{2}-y^{2}$, el descenso de gradiente con paso $\eta$ actúa por coordenadas: $x\leftarrow(1-2\eta)x$ e $y\leftarrow(1+2\eta)y$. La coordenada $x$ se contrae y la $y$ se expande. Si el punto inicial tiene $y_0=0$, el método converge a la silla; cualquier perturbación $y_0\ne 0$ crece geométricamente y escapa de ella (\cref{eje:1-silla}).
\end{ejemplo}

\begin{observacion}[El paisaje de las redes profundas]\label{obs:paisaje}
Se conocen varios hechos empíricos sobre la pérdida de las redes profundas. (i) En alta dimensión, los puntos críticos de pérdida alta son mayoritariamente sillas, no mínimos locales: un mínimo exige que todas las curvaturas sean positivas a la vez~\cite{dauphin}, \cite[§8.2.3]{gbc}. (ii) El espectro del hessiano se concentra cerca de cero, con un grupo menor de autovalores alejados que dependen de los datos~\cite{sagun}; el problema está muy mal condicionado~\cite[§8.2.1]{gbc}. (iii) Los lotes muy grandes tienden a llevar a mínimos más agudos, que se han asociado a una peor generalización~\cite{keskar}. (iv) Las conexiones residuales suavizan visiblemente el paisaje~\cite{li-paisaje}; se estudian en el módulo~11. Los temas siguientes desarrollan algoritmos que funcionan razonablemente en este entorno aunque su análisis se haga en el caso convexo.
\end{observacion}

\begin{lecturas}
\item Boyd y Vandenberghe~\cite{bv}, cap.~2 «Convex sets», §3.1 «Basic properties and examples» y §9.1 «Unconstrained minimization problems».
\item Stanford EE364a~\cite{ee364a}, transparencias del curso correspondientes a los capítulos~2, 3 y~9 de~\cite{bv}.
\item Goodfellow, Bengio y Courville~\cite{gbc}, §8.1 «How Learning Differs from Pure Optimization» y §8.2 «Challenges in Neural Network Optimization».
\item Zhang \emph{et al.}~\cite{d2l}, §12.1 «Optimization and Deep Learning» y §12.2 «Convexity».
\end{lecturas}

\section*{Ejercicios del tema 1}
\addcontentsline{toc}{section}{Ejercicios del tema 1}

\begin{ejercicio}{1}{eje:1-cuartica}
Sea $f(x)=x^{4}$ en $\R$. Justifique si es (a) convexa, (b) fuertemente convexa, (c) $L$-suave en todo $\R$ para algún $L$. ¿Contradice el apartado (b) que tenga un único minimizador?
\end{ejercicio}

\begin{ejercicio}{1}{eje:1-silla}
Clasifique el punto $(0,0)$ de $f(x,y)=x^{2}-y^{2}$ y explique por qué el descenso de gradiente, empezando exactamente en $(1,0)$, converge a él. ¿Qué ocurre si empieza en $(1,\delta)$ con $\delta\ne 0$? Tome $\eta=0{,}1$ y $\delta=10^{-6}$ y calcule $y_{50}$.
\end{ejercicio}

\begin{ejercicio}{2}{eje:1-cuadratica}
Halle el minimizador y el valor mínimo de $h(x)=\tfrac12x\T Qx+x\T c+b$ con
$Q=\begin{pmatrix}2&1\\1&2\end{pmatrix}$, $c=(-3,0)\T$ y $b=1$.
\fuente{adaptado de Zhang \emph{et al.}~\cite{d2l}, §12.6, ejercicio~3.}
\end{ejercicio}

\begin{ejercicio}{2}{eje:1-constantes}
Sea $f(x,y)=x^{2}+2xy+5y^{2}$. Calcule su hessiano, demuestre que $f$ es fuertemente convexa y obtenga $\mu$, $L$ y $\kappa$.
\end{ejercicio}

\begin{ejercicio}{3}{eje:1-maximo}
Sean $f$ y $g$ funciones convexas en $\R^{n}$. (a) Demuestre que $h(x)=\max\{f(x),g(x)\}$ es convexa. (b) Muestre con un contraejemplo que $\min\{f(x),g(x)\}$ no es convexa en general.
\fuente{adaptado de Zhang \emph{et al.}~\cite{d2l}, §12.2, ejercicio~3.}
\end{ejercicio}

\begin{ejercicio}{3}{eje:1-fuerte}
Sea $f$ diferenciable y $\mu$-fuertemente convexa, con un minimizador $\theta^{\ast}$. Demuestre que (a) $f(\theta)-f^{\ast}\ge\frac{\mu}{2}\norm{\theta-\theta^{\ast}}^{2}$ para todo $\theta$; (b) el minimizador es único; (c) $\norm{\nabla f(\theta)}\ge\mu\norm{\theta-\theta^{\ast}}$.
\end{ejercicio}
